% 图契约：回旋加速器的半圆轨道与窄缝；图中的端口、几何、箭头和方程对应本题。
% 版面：沿用前三模块的几何图+推导分区；目标印宽17cm，简图至中等密度。
% 布局：左侧真实结构，右侧关键关系；真实电路边不作为架构流程箭头处理。
\documentclass[tikz,border=7pt]{standalone}
\usepackage{fontspec,xeCJK,amsmath,amssymb}
\IfFontExistsTF{Times New Roman}{\setmainfont{Times New Roman}}{\setmainfont{TeX Gyre Termes}}
\IfFontExistsTF{Heiti SC}{\setCJKmainfont[AutoFakeBold=1.5]{Heiti SC}}{\setCJKmainfont[AutoFakeBold=1.5]{FandolHei-Regular}}
\usetikzlibrary{arrows.meta,bending,calc,patterns,positioning,decorations.pathmorphing}
\definecolor{ink}{HTML}{243745}
\definecolor{blue}{HTML}{275DAB}
\definecolor{green}{HTML}{177C63}
\definecolor{red}{HTML}{BD4935}
\definecolor{orange}{HTML}{AA6C16}
\definecolor{purple}{HTML}{715393}
\tikzset{every node/.style={font=\fontsize{9.3}{12.5}\selectfont,text=ink},
arrow/.style={-{Stealth[length=5pt 1.5,width'=0pt 3]},shorten >=2pt,shorten <=1pt,line width=1pt},
arrow short/.style={-{Stealth[length=3pt 1.0,width'=0pt 2.5]},shorten >=1pt,shorten <=0pt,line width=.8pt},
axis/.style={arrow,draw=ink!60,line width=.6pt},
vec/.style={arrow,draw=blue},
wire/.style={draw=ink!75,line width=.8pt},
guide/.style={dashed,draw=ink!45,line width=.5pt},
card/.style={rounded corners=2pt,draw=blue!65,fill=blue!4,align=center,text width=6.5cm,inner xsep=9pt,inner ysep=7pt},
title/.style={font=\fontsize{10.5}{14}\selectfont\bfseries,text width=15cm,align=center},
label/.style={font=\fontsize{8.4}{11}\selectfont,align=center},
note/.style={label,text width=6.1cm}}
\begin{document}
\begin{tikzpicture}[x=1cm,y=1cm]
\path[use as bounding box] (-8,-3.5) rectangle(8,3.65);
\node[title] at (0,3.25) {理论图\quad 回旋加速器的半圆轨道与窄缝};
\begin{scope}[shift={(-4.05,0)}]\filldraw[fill=blue!3,draw=ink!50](-.12,-2)arc(-90:-270:2)--cycle;\filldraw[fill=green!3,draw=ink!50](.12,-2)arc(-90:90:2)--cycle;\draw[orange,thick](0.12,0)arc(-90:90:0.48);\draw[arrow short,orange](0.6,0.48)arc(0:25:0.48);\draw[red,thick](0.12,0.96)--(-0.12,0.96);\draw[orange,thick](-0.12,0.96)arc(90:270:0.6788225099390857);\draw[arrow short,orange](-0.7988225099390857,0.2811774900609143)arc(180:205:0.6788225099390857);\draw[red,thick](-0.12,-0.39764501987817136)--(0.12,-0.39764501987817136);\draw[orange,thick](0.12,-0.39764501987817136)arc(-90:90:0.831384387633061);\draw[arrow short,orange](0.951384387633061,0.43373936775488964)arc(0:25:0.831384387633061);\draw[red,thick](0.12,1.2651237553879506)--(-0.12,1.2651237553879506);\draw[orange,thick](-0.12,1.2651237553879506)arc(90:270:0.96);\draw[arrow short,orange](-1.08,0.30512375538795067)arc(180:205:0.96);\draw[<->,red](-.12,2.25)--(.12,2.25);\node[above]at(0,2.3){交变电压 $U$};\node[label]at(0,-2.45){$q>0,\ B$ 向纸内；窄缝增能 $qU$};\end{scope}\node[card,draw=blue!65,fill=blue!4]at(3.9,1.8){$\omega_c=|q|B/m$};\node[card,draw=green!65,fill=green!4]at(3.9,0){$\Delta K=|q|U$（每次过缝）};\node[card,draw=orange!65,fill=orange!4]at(3.9,-1.8){$R_N^2=2mN U/(|q|B^2)$};
\end{tikzpicture}\end{document}